% Separate post-CP26 proved geometric comparison. No finite-family effectivity or hereditary projective count.
\subsection{Bounded-period root anchors and canonical full orders}

\begin{proposition}[Proved: a stable twisted anchor and its induced lattices]
\label{D:stable-root-anchor-canonical-comparison}
Use the fixed complete surface base, a degree-$t$ normal punctured
cover $A/R$, and its original comparable-cycle model
$Y_0\to\operatorname{Spec}A$ from
Proposition~\ref{D:uniform-comparable-reduction-cycle}.
Let $\rho:Y\to Y_0$ be a regular point-blowup refinement, and
write $C=\rho^*C_0$, where
$(x,y)\mathcal O_{Y_0}=\mathcal O_{Y_0}(-C_0)$ with $x,y\in J$.
All full bundles below mean the actual pullbacks from $Y_0$;
they are not newly canonicalized on $Y$.

Suppose the reduced exceptional support on $Y$ is SNC and
$\mathcal X\to Y$ is a tame root stack on that support. Let
$\mathcal G\to\mathcal X$ be a banded $\mu_n$-gerbe,
$(n,p)=1$, extending a realized coefficient class on the
puncture $U$. Fix a scalar weight and a rank-$h$ twisted
bundle $\mathcal A_U$ which admits a vector-bundle extension
to $\mathcal G$. Put
\[
 P_0=j_*\mathcal End(\mathcal A_U),\qquad
 P=\rho^*\operatorname{full}_{Y_0}(P_0),
 \qquad j:U\hookrightarrow\operatorname{Spec}A.
                                                        \tag{D.BA.1}
\]
Descent of scalar-weight-zero bundles to $U$ is understood.
There exists a vector-bundle extension $\mathcal A_1$ of
$\mathcal A_U$, of the same rank $h$, with an actual left
$P$-action. If $q:\mathcal G\to Y$ is the coarse map, then
$\mathcal B_1=q_*\mathcal End(\mathcal A_1)$ is locally free
over $Y$ and, inside its fixed generic endomorphism algebra,
\[
                  P\subset\mathcal B_1\subset P^\vee
                             \subset P(C).             \tag{D.BA.2}
\]
The middle order need not be Azumaya on the exceptional divisor.

For every punctured twisted bundle $\mathcal F_U$ of the
same scalar weight, set
\[
 \begin{aligned}
 Q_0&=j_*\mathcal Hom(\mathcal A_U,\mathcal F_U),&
 R_0&=j_*\mathcal Hom(\mathcal F_U,\mathcal A_U),\\
 Q&=\rho^*\operatorname{full}_{Y_0}(Q_0),&
 R_1&=\rho^*\operatorname{full}_{Y_0}(R_0).
 \end{aligned}                                           \tag{D.BA.3}
\]
There is a vector-bundle extension $\mathcal F$ on
$\mathcal G$, of the original rank, for which the ordinary
locally free coarse Hom bundle satisfies
\[
 Q\subset q_*\mathcal Hom(\mathcal A_1,\mathcal F)
       \subset R_1^\vee\subset Q(C).                  \tag{D.BA.4}
\]
Every identification and inclusion restricts to the specified
one on $U$. These are assertions about individual geometric
bundles and actual sheaf lattices, not an assertion about
arbitrary relative reflexive hulls.
\end{proposition}

\begin{proof}
The canonical order and module actions are those of
Proposition~\ref{D:canonical-order-Morita-conductor}.
We recall the dual comparison to fix its precise scope.
For any finite reflexive module $L_0$, put
$L=\operatorname{full}_{Y_0}(L_0)$ and
$L'=\operatorname{full}_{Y_0}(L_0^\vee)$.
The natural evaluation gives $L'\subset L^\vee$.
For $z\in J$, the first-power Noether factorization writes
$z\operatorname{id}_{L_0}=vu$ through a finite free module.
The entries $u_i$ are sections of $L'$; hence for a local
$\lambda\in L^\vee$,
\[
                 z\lambda=\sum_i\lambda(v_i)u_i\in L'.
                                                        \tag{D.BA.5}
\]
Using $z=x,y$ proves $L^\vee(-C_0)\subset L'\subset L^\vee$.
Pull back these actual maps to $Y$.
The trace pairings on the puncture extend reflexively over
the normal surface, so $P_0=P_0^\vee$ and $R_0=Q_0^\vee$.
Thus, on the refined model as well,
\[
 P\subset P^\vee\subset P(C),\qquad
 R_1\subset Q^\vee\subset R_1(C),\qquad
 Q\subset R_1^\vee\subset Q(C).                       \tag{D.BA.6}
\]
Neither a boundary Azumaya property nor a division by $h$
is used in these dual inclusions.

Start with any vector-bundle extension $\mathcal A^{(0)}$
of $\mathcal A_U$. Form the coherent image
\[
 \mathcal L=\operatorname{im}
   \bigl(q^*P\otimes\mathcal A^{(0)}
                   \longrightarrow\mathcal A^{(0)}(*E)\bigr),
             \qquad \mathcal A_1=\mathcal L^{\vee\vee}.
                                                        \tag{D.BA.7}
\]
The map is the prescribed generic action. Coherence of the
image follows locally by clearing finitely many denominators;
on a finite cover they can be cleared by some exceptional
Cartier multiple. No uniform bound for that preliminary
denominator is needed or claimed. The unit makes
$\mathcal A^{(0)}\subset\mathcal L$. Associativity of $P$
makes $\mathcal L$ stable under $q^*P$.

The gerbe and the root stack have regular two-dimensional
etale scheme charts. On these charts a coherent reflexive
sheaf is locally free, and maps to a reflexive target extend
uniquely once they extend at height one
\citep[Tags 0AVC and 0B3N; Lemmas 31.13.12--31.13.15]
{C_StacksRootAnchorReflexive}.
Reflexive hull and these maps commute with the etale
transition maps, so the construction descends to
$\mathcal G$. It gives a rank-$h$ vector bundle and an
actual $P$-action on $\mathcal A_1$. The algebra identities
hold everywhere because they hold on the dense puncture.

Scalar-weight-zero Hom descends from the gerbe to the root
stack. Its coarse pushforward is locally free: on a root
chart the coordinate algebra is finite locally free over
the regular coarse chart, a bundle is projective over that
algebra, and taking invariants under the finite tame root
group is a direct summand. Thus its invariants are projective
over the coarse chart. This also proves coherence and the
local-freeness assertions for the later coarse Hom bundles.

The action embeds $P$ into $\mathcal B_1$. If $b$ is a
section of $\mathcal B_1$ and $a$ a section of $P$, the
trace of their product on $\mathcal A_1$ is regular on
$\mathcal G$ and descends to $\mathcal O_Y$.
It agrees on $U$ with the ordinary matrix trace pairing
used in (D.BA.6). Consequently $b\in P^\vee$.
This proves (D.BA.2), including at the closed points.

For (D.BA.4), use the right $P$-module $Q$ to define
\[
 \mathcal F=
   \bigl((q^*Q\otimes_{q^*P}\mathcal A_1)_{
                  \mathrm{torsion\,free}}\bigr)^{\vee\vee}.
                                                        \tag{D.BA.8}
\]
Here torsion-free quotient means the image in the specified
generic twisted vector space. On $U$, ordinary matrix
Morita evaluation identifies that space with $\mathcal F_U$.
The regular-chart argument just given makes (D.BA.8) a
vector bundle of its original rank. Tensor evaluation
gives the actual map
$Q\to q_*\mathcal Hom(\mathcal A_1,\mathcal F)$.
It is injective because it is an isomorphism generically.

The canonical composition $R_1\otimes Q\to P$ also gives
\[
 q^*R_1\otimes
   (q^*Q\otimes_{q^*P}\mathcal A_1)
                 \longrightarrow\mathcal A_1.
\]
It kills torsion. At height one the torsion-free quotient
is already reflexive; hence this map extends to (D.BA.8),
by the same height-one extension criterion. It provides
$R_1\to q_*\mathcal Hom(\mathcal F,\mathcal A_1)$.
Pair this map with $q_*\mathcal Hom(\mathcal A_1,\mathcal F)$
by composition and trace on $\mathcal A_1$. The pairing is
regular and is the perfect generic pairing $Q_0\times R_0$.
Therefore the coarse Hom lattice lies in $R_1^\vee$.
The last inclusion in (D.BA.4) is (D.BA.6).
\end{proof}

\begin{corollary}[Proved: bounded-rank generating anchors after a bipartite refinement]
\label{D:bounded-period-generating-root-anchor}
In the preceding cover setup, suppose the geometric Brauer
period is $M\le B$, with $(M,p)=1$. After an $O_R(t)$
point-blowup SNC refinement and one further blowup of each
SNC node, use the common $M$-root stack of the entire
exceptional divisor. The actual coefficient gerbe admits
a scalar-weight-one generating bundle $\mathcal A_{\rm gen}$
of rank at most $B^4$.

More precisely, choose a rank-$h$ punctured anchor with
$h=M^2$ by the qualitative residual-anchor construction
in Proposition~\ref{D:SNC-general-weighted-presentation},
and choose $\mathcal A_1$ as above. There are $d=M^2$
line twists, each canonically trivial on $U$, such that
their sum tensored with $\mathcal A_1$ is a generating
bundle. Put $P^{[d]}=\operatorname{Mat}_d(P)$ and
$Q^{[d]}=Q^{\oplus d}$. Under their canonical punctured
identifications, its actual ordinary coarse order satisfies
\[
 P^{[d]}(-C)\subset
 q_*\mathcal End(\mathcal A_{\rm gen})
                  \subset P^{[d]}(2C),                \tag{D.BA.9}
\]
and the induced extension $\mathcal F$ of (D.BA.8) satisfies
\[
 Q^{[d]}(-C)\subset
 q_*\mathcal Hom(\mathcal A_{\rm gen},\mathcal F)
                  \subset Q^{[d]}(C).                 \tag{D.BA.10}
\]
All pole constants are independent of the covering degree,
coefficient order, rank of $\mathcal F_U$, and number of
exceptional components. They remain valid for any scalar
weight prime to the coefficient order.
\end{corollary}

\begin{proof}
Lemma~\ref{D:linear-SNC-refinement-centres} and
Corollary~\ref{D:SNC-refinement-inherited-cycles} give
the first SNC refinement with $O_R(t)$ nodes and the
inherited Cartier cycles. Blow up each of its nodes once.
All original components are now pairwise disjoint, and all
new exceptional components are pairwise disjoint. Each node
joins one component of each kind. This is an actual bipartite
divisor, also when the old graph had cycles or parallel edges.
These are $O_R(t)$ extra multiplicity-two blowups; all inherited
cycle and cohomology identities still hold. Write the two
sets of components as $I_0,I_1$.

The actual coefficient class extends to its mixed residue-root
stack by Lemma~\ref{D:SNC-general-Brauer-geometry}. Its residue
orders divide $M$, so pull it to the common $M$-root stack.
The geometric-period-$M$ residual anchor of rank $M^2$,
followed by formal lifting and algebraization, is precisely
the qualitative existence construction (D.GP.3)--(D.GP.5).
Use its restriction as $\mathcal A_U$. No finite etale
splitter of degree bounded in $M$ is inferred.

Let $\mathcal D_i$ be the $M$-root divisor over $E_i$, and put
\[
 \mathcal T_{a,b}=\mathcal O_{\mathcal X}
       \left(a\sum_{i\in I_0}\mathcal D_i
                    +b\sum_{i\in I_1}\mathcal D_i\right),
 \quad 0\le a,b<M,\qquad
 \mathcal A_{\rm gen}
       =\bigoplus_{a,b}\mathcal A_1\otimes\mathcal T_{a,b}.
                                                        \tag{D.BA.11}
\]
Their effective rational divisor coefficients on $Y$ are
$a/M$ or $b/M$, all less than one. Every coefficient of
$C$ along the exceptional support is a positive integer.
The canonical root sections therefore give regular maps
\[
 \mathcal O\longrightarrow\mathcal T_{a,b}
                      \longrightarrow\mathcal O(C),
 \qquad
 \mathcal O(-C)\longrightarrow
       \mathcal T_{a',b'}\mathcal T_{a,b}^{-1}
                      \longrightarrow\mathcal O(C),   \tag{D.BA.12}
\]
all restricting to the chosen identity on $U$.

At a node, the twists in (D.BA.11) give every character of
$\mu_M\times\mu_M$ once; at a one-component point they give
every character of $\mu_M$. This still generates in the
presence of a noncommuting central gerbe inertia. Indeed let
$D$ be that finite tame inertia group, $Z$ its scalar centre,
and $H=D/Z$ its abelian root quotient. For any nonzero
representation $V$ of the specified scalar weight,
\[
 V\otimes k[H]\simeq
       \operatorname{Ind}_Z^D\operatorname{Res}_Z^D V.
                                                        \tag{D.BA.13}
\]
For every irreducible $W$ of that scalar weight, Frobenius
reciprocity gives
$\operatorname{Hom}_D(W,V\otimes k[H])
 =\operatorname{Hom}_Z(W,V)\ne0$.
Semisimplicity holds because all inertia orders are prime to
$p$. Hence every scalar-weight-one inertia simple occurs.
Away from the exceptional divisor the scalar centre alone
remains, and a nonzero bundle suffices. On the finite root
charts, this fibrewise statement and Nakayama give a projective
generator, exactly as in the generating-bundle proof of
Lemma~\ref{D:SNC-general-Brauer-geometry}.
Thus (D.BA.11) is an actual generating bundle of rank
$hM^2=M^4\le B^4$.

For each matrix block in its endomorphism algebra,
(D.BA.12), exact tame pushforward, and (D.BA.2) give
\[
 P(-C)\subset\mathcal B_1(-C)
 \subset q_*\mathcal Hom(
       \mathcal A_1\mathcal T_{a,b},
       \mathcal A_1\mathcal T_{a',b'})
 \subset\mathcal B_1(C)\subset P(2C).
                                                        \tag{D.BA.14}
\]
Summing the blocks proves (D.BA.9). In particular the claim
is not $P^{[d]}\subset q_*\mathcal End(\mathcal A_{\rm gen})$:
matrix units between unequal inertia weights can require
one Cartier factor.

Put $E_1=q_*\mathcal Hom(\mathcal A_1,\mathcal F)$.
The positivity in the first part of (D.BA.12) gives
\[
 E_1(-C)\subset
 q_*\mathcal Hom(\mathcal A_1\mathcal T_{a,b},\mathcal F)
                  \subset E_1.                        \tag{D.BA.15}
\]
Combine this with (D.BA.4) and sum all $M^2$ blocks to
obtain (D.BA.10). All constructions and trace identities
work for any fixed scalar weight coprime to $n$, giving
the last assertion.
\end{proof}

\begin{remark}[What the comparison does not establish]
The canonical bundles $P,Q$ are generated after the inherited
strong-model twist, and have the previously proved degree
budgets. The intermediate coarse lattices in (D.BA.9)--(D.BA.10)
are not thereby proved generated or cohomologically bounded:
a divisor sandwich alone does not bound their instability.
Nor has this construction made the arbitrary choices of
reflexive hulls flat in Artinian families. The exact fixed
generating-bundle Morita equivalence is available once the
anchor has been chosen, but arbitrary finite nilpotent order
actions are not asserted to produce locally free root-stack
lattices. Simultaneous Frobenius comparison, finite-data
effectivity, and a degree-linear whole-cycle presentation
remain separate requirements. These two results are geometric
comparison inputs, not a hereditary projective count or an
independent review movement.
\end{remark}
